\chapter{Trotter-based Hamiltonian Simulation}\label{ch:sec04}

Given an initial state $\ket{\psi_0}$ and a Hamiltonian $H$, we evaluate
$\ket{\psi(t)} = \ee^{-\ii H t}\ket{\psi_0}$, i.e.\ we solve the ODE
\begin{equation}
  \partial_t \psi(t) = -\ii H\psi(t), \qquad \psi(0) = \psi_0.
\end{equation}

\section{Trotter splitting}\label{sec:sec04-trotter}

Consider $H = H_1 + H_2$, where $\ee^{-\ii H_1 t}$ and $\ee^{-\ii H_2 t}$ can be
computed efficiently, but in general $[H_1, H_2] \neq 0$. The Lie--Trotter
product formula gives
\begin{equation}
  \ee^{-\ii H t} = \lim_{L \to \infty}\bigl(\ee^{-\ii \frac{t}{L} H_1}\ee^{-\ii \frac{t}{L} H_2}\bigr)^L.
\end{equation}
Taking $L$ as a finite number and $\Delta t = \tfrac{t}{L}$ gives the first-order
Trotter method, with local error
\begin{equation}
  \norm{\ee^{-\ii \Delta t H} - \ee^{-\ii \Delta t H_1}\ee^{-\ii \Delta t H_2}} = \bigo(\Delta t^2).
\end{equation}
Therefore, to perform Hamiltonian simulation up to time $t$,
\begin{equation}
  \norm{\ee^{-\ii t H} - \bigl(\ee^{-\ii \frac{t}{L} H_1}\ee^{-\ii \frac{t}{L} H_2}\bigr)^L}
  = \bigo(\Delta t^2 \cdot L) = \bigo\!\Bigl(\frac{t^2}{L}\Bigr),
\end{equation}
so to reach precision $\varepsilon$ in operator norm we need
$L = \bigo(t^2 \varepsilon^{-1})$. This can be improved by the second-order
Trotter method (symmetric Trotter splitting, or Strang splitting),
\begin{equation}
  \norm{\ee^{-\ii \Delta t H} - \ee^{-\ii \frac{\Delta t}{2} H_2}\ee^{-\ii \Delta t H_1}\ee^{-\ii \frac{\Delta t}{2} H_2}} = \bigo(\Delta t^3),
\end{equation}
so to achieve precision $\varepsilon$ we need
$L = \bigo(t^{3/2}\varepsilon^{-1/2})$. For a higher-order Trotter method, such as
the $2p$-th order Suzuki formula with local error $(\Delta t)^{2p+1}$, we need
\begin{equation}
  L = \bigo\!\bigl(t^{\frac{2p+1}{2p}}\varepsilon^{-\frac{1}{2p}}\bigr).
\end{equation}

\begin{example}[Transverse-field Ising model]
For the one-dimensional transverse-field Ising model (TFIM) with nearest-neighbor
interaction,
\begin{equation}
  H = \underbrace{-\sum_{i=1}^{n-1} Z_i Z_{i+1}}_{H_1} \underbrace{-\, g\sum_{i=1}^{n} X_i}_{H_2},
  \qquad Z = \begin{psmallmatrix} 1 & 0 \\ 0 & -1 \end{psmallmatrix}.
\end{equation}
Then
\begin{equation}
  \ee^{-\ii t H_1} = \ee^{\ii t \sum_{i=1}^{n-1} Z_i Z_{i+1}} = \prod_{i=1}^{n-1} \ee^{\ii t Z_i Z_{i+1}}, \qquad
  \ee^{-\ii t H_2} = \ee^{\ii t g \sum_{i=1}^{n} X_i} = \prod_{i=1}^{n} \ee^{\ii t g X_i},
\end{equation}
where each rotation only involves the qubits $i, j$.
\end{example}

\begin{example}[Particle in a potential]
Let $H = -\Delta_r + V(r) = H_1 + H_2$, with $r \in \Omega = [0,1]^d$ and periodic
boundary conditions. Then $\ee^{-\ii H_1 t}$ can be efficiently performed since
$H_1$ is diagonalizable in Fourier space, and $\ee^{-\ii H_2 t}$ can be
efficiently performed since $V(r)$ is diagonal in real space.
\end{example}

\section{Commutator-type error bound}\label{sec:sec04-commutator}

We refine the previous error estimate. Let
$\tilde{u}(t) = \ee^{-\ii H_1 t}\ee^{-\ii H_2 t}$. Then
\begin{align}
  \ii\partial_t \tilde{u}(t)
  &= H_1 \ee^{-\ii H_1 t}\ee^{-\ii H_2 t} + \ee^{-\ii H_1 t} H_2 \ee^{-\ii H_2 t} \notag\\
  &= (H_1 + H_2)\ee^{-\ii H_1 t}\ee^{-\ii H_2 t} - H_2 \ee^{-\ii H_1 t}\ee^{-\ii H_2 t} + \ee^{-\ii H_1 t} H_2 \ee^{-\ii H_2 t} \notag\\
  &= H\tilde{u}(t) + [\ee^{-\ii H_1 t}, H_2]\ee^{-\ii H_2 t},
\end{align}
with $\tilde{u}(0) = I$. Let $u(t) = \ee^{-\ii H t}$. By Duhamel's principle,
\begin{equation}
  \tilde{u}(t) = u(t) - \ii\int_0^t \ee^{-\ii H(t-s)}[\ee^{-\ii H_1 s}, H_2]\ee^{-\ii H_2 s}\,\dd s,
\end{equation}
so
\begin{equation}
  \norm{\tilde{u}(t) - u(t)} \le \int_0^t \norm{[\ee^{-\ii H_1 s}, H_2]}\,\dd s.
\end{equation}
Let $G(t) = [\ee^{-\ii H_1 t}, H_2]\ee^{\ii H_1 t} = \ee^{-\ii H_1 t} H_2 \ee^{\ii H_1 t} - H_2$,
so that $G(0) = 0$ and $\ii\partial_t G(t) = \ee^{-\ii H_1 t}[H_1, H_2]\ee^{\ii H_1 t}$.
Hence $\norm{[\ee^{-\ii H_1 t}, H_2]} = \norm{G(t)} \le t\norm{[H_1, H_2]}$, and
therefore
\begin{equation}
  \norm{\tilde{u}(t) - u(t)} \le \int_0^t s\norm{[H_1, H_2]}\,\dd s
  \le \frac{t^2}{2}\norm{[H_1, H_2]} \le t^2 \nu^2,
\end{equation}
and
\begin{equation}
  \norm{\ee^{-\ii \Delta t H_1}\ee^{-\ii \Delta t H_2} - \ee^{-\ii \Delta t H}}
  \le \frac{\Delta t^2}{2}\norm{[H_1, H_2]} \le (\Delta t)^2 \nu^2,
\end{equation}
with $\nu = \max\{\norm{H_1}, \norm{H_2}\}$. The bound involving
$\norm{[H_1,H_2]}$ is the \emph{commutator-norm} error estimate, and the one
involving $\nu$ is the \emph{operator-norm} error estimate.

\begin{example}
For $H = -\sum_i Z_i Z_{i+1} - g\sum_i X_i$ (with $H_1, H_2$ as above),
$\norm{H_1}, \norm{H_2} = \bigo(n)$, so $\nu^2 = \bigo(n^2)$; while
$[Z_i Z_j, X_k] \neq 0$ only if $k = i$ or $k = j$, so
$\norm{[H_1, H_2]} = \bigo(n)$. To reach accuracy $\varepsilon$: the
operator-norm bound gives $L \sim \bigo(n^2/\varepsilon)$, while the
commutator-norm bound gives $L \sim \bigo(n/\varepsilon)$.
\end{example}

\begin{example}
For $H = -\Delta + V$ (particle in a potential), take $d = 1$, $\Omega = [0,1]$
discretized using a uniform grid of size $N$, with $V$ smooth and bounded. Then
$\norm{H_1} = \bigo(N^2)$, $\norm{V} = \bigo(1)$, so $\nu^2 = \bigo(N^4)$, giving
$L = \bigo(N^2/\varepsilon)$. On the other hand, for any smooth function
$\psi(r)$,
\begin{equation}
  [H_1, H_2]\psi = \Bigl[-\frac{\dd^2}{\dd r^2}, V\Bigr]\psi = -V''\psi - 2V'\psi',
\end{equation}
so $\norm{[H_1, H_2]\psi} \le \norm{V''} + 2\norm{V'}\cdot\norm{\psi'} = \bigo(N)$,
where the three terms are $\bigo(1)$, $\bigo(1)$, and $\bigo(N)$ respectively.
\end{example}
